% !TeX root=../main.tex
\chapter{تحلیل ناهمانندی در فضای بردارهای توزیع احتمال}
\label{app:probability-dissimilarity}
در این پیوست، سنجه‌های اطلاعاتیِ مقایسهٔ توزیع‌های احتمال، نمایش هیلبرتی واگرایی جنسن--شنون، متریک و شباهت حاصل از آن و خواص هندسی این فضا را به‌صورت یکپارچه بررسی می‌کنیم. اثبات‌های تفصیلی پس از معرفی مفاهیم آمده‌اند. کاربرد این سنجه‌ها در پیش‌بینی برچسب و تحلیل خطای تصمیم در فصل~\ref{cmp:ch:method} بررسی می‌شود.

\section{ناهمانندی، فاصله و شباهت توزیع‌ها}
\label{cmp:subsec:ch3-probability-distance-similarity}
\label{cmp:subsec:ch3-probability-dissimilarity-families}
توزیع‌های برچسب همسایگی و نماینده‌های طبقات،
اعضای فضای $\Delta_K$ هستند که در
بخش~\ref{cmp:subsec:ch3-class-centers}
معرفی شد. این فضا زیرمجموعه‌ای از
$\mathbb R^K$ است، اما مؤلفه‌های بردارهای آن
نامنفی و دارای مجموع یک‌اند و ممکن است صفر باشند. برای سنجش نزدیکی
دو توزیع و مقایسه نمونه با مرکز،
سنجه‌ای لازم است که با این تعبیر احتمالاتی
سازگار باشد. البته فاصله اقلیدسی نیز در این فضا
قابل تعریف است؛ ما سنجه‌های اطلاعاتی را
به دلیل معنای احتمالاتی آن‌ها بررسی می‌کنیم.
\emph{واگرایی} را نباید با \emph{متریک فاصله}
یا \emph{تابع شباهت} یکسان دانست:
یک متریک علاوه بر نامنفی‌بودن و صفرشدن
فقط برای دو توزیع برابر، تقارن و
نامساوی مثلثی را نیز ارضا می‌کند.
تابع شباهت، برخلاف سنجهٔ ناهمانندی،
مقدار بزرگ‌تر را نشانهٔ سازگاری بیشتر می‌گیرد.
هر سنجه باید با توجه به جهت، دامنهٔ مقدارها،
رفتار مؤلفه‌های صفر و نقش آن در تصمیم انتخاب شود.
در ادامه، نخست محدودیت‌های KL را می‌گوییم،
سپس با یک نمایش هیلبرتی به واگرایی جنسن--شنون،
متریک جذر آن و شباهت متناظر می‌رسیم.
این واگرایی پیش‌تر در نظریهٔ اطلاعات معرفی شده است
\cite{Lin1991JensenShannon} و نمایش هیلبرتی آن را
در ادامه با استدلال صریح بررسی می‌کنیم.
انتخاب سنجه، تعریف مرکز میانگین را تغییر نمی‌دهد.

\paragraph{واگرایی KL و محدودیت‌های آن.}
واگرایی KL نخست از تجزیهٔ لگاریتم درست‌نمایی
در بخش~\ref{cmp:subsec:ch3-mean-center-justification}
ظاهر شد و صورت ریاضی آن در
رابطه~\eqref{cmp:eq:ch3-KL-definition} آمده است.
در اینجا آن را به‌عنوان یکی از انتخاب‌های ممکن
برای \emph{مقایسهٔ بازنمایی نمونه و مرکز} بررسی می‌کنیم؛
نقش آن در توجیه مرکز میانگین با انتخاب فاصلهٔ تصمیم
یکسان نیست.
این سنجه نامنفی است و فقط برای $P=Q$ صفر می‌شود،
اما جهت‌دار و عموماً نامتقارن است،
نامساوی مثلثی را ندارد و با وجود
$p_j>0$ و $q_j=0$ می‌تواند بی‌نهایت شود.
پس KL با وجود تعبیر اطلاعاتی مهم آن،
متریک فاصله محسوب نمی‌شود.

\paragraph{انگیزهٔ نمایش هیلبرتی.}
محدودیت‌های KL ما را به بررسی ناهمانندی‌ای
متقارن، کراندار و قابل‌تعریف حتی در حضور
مؤلفه‌های صفر هدایت می‌کنند. بااین‌حال،
صرفِ متقارن‌بودن یک واگرایی برای داشتن
نامساوی مثلثی کافی نیست. راه مناسب آن است
که توزیع‌ها را به فضایی هیلبرتی ببریم و
فاصلهٔ اقلیدسیِ تعمیم‌یافتهٔ تصاویر را
به‌عنوان متریک در نظر بگیریم.
برای حفظ تعبیر شباهت، مطلوب است تصاویر
\emph{نرم ثابت} نیز داشته باشند؛ در این صورت
ضرب داخلی، پس از نرمال‌سازی، شباهتی سازگار
با فاصلهٔ تصاویر به دست می‌دهد.

نگاشت زیر انتخابی دلخواه از میان نگاشت‌های
هیلبرتی نیست: شکل نماییِ مؤلفه‌های آن با
بازنمایی انتگرالی عبارت‌های لگاریتمی سازگار
است و عامل مقیاس به‌گونه‌ای برگزیده می‌شود
که نرم مربع هر مؤلفه با احتمال همان مؤلفه
متناسب باشد. بنابراین قید مجموع یکِ
توزیع‌ها، نرم کل تصویر را ثابت می‌کند.
پس از ساخت نگاشت، نشان می‌دهیم مربع
فاصلهٔ تصاویر دقیقاً برابر JSD است؛
در نتیجه جذر JSD متریک و شباهت نرمال‌شده
دارای تفسیر ضرب داخلی خواهد بود.
جزئیات هویت انتگرالی و شرایط مرزی
در بخش~\ref{app:probability-dissimilarity-proofs} آمده‌اند.

\paragraph{نگاشت توزیع‌ها به فضای هیلبرت.}
فضای هیلبرت $\mathcal H=L^2((0,\infty);\mathbb R^K)$ را در نظر گرفته و هر توزیع $P=(p_1,\ldots,p_K)$ را به تابع برداری زیر نگاشت می‌کنیم:
\begin{equation}\label{cmp:eq:ch3-js-Hilbert-map}
\boxed{
\Phi(P)(t)=
\bigl(\phi_{p_1}(t),\ldots,\phi_{p_K}(t)\bigr),
\qquad
\phi_a(t)=\frac{1-e^{-at}}{2t},\quad t>0.
}
\end{equation}
این نگاشت تعریف‌شده و یک‌به‌یک است و خواص زیر را دارد: برای هر $t>0$، مؤلفه $\phi_a(t)$ نسبت به احتمال $a$ اکیداً صعودی و مقعر است؛ برای هر $a>0$، نسبت به $t$ نزولی بوده و بیشینه حدی آن در $t\to0^+$ برابر $a/2$ است. بنابراین نگاشت، ترتیب مؤلفه‌های احتمال را حفظ می‌کند و دو توزیع متفاوت را به دو تصویر متفاوت می‌برد. افزون بر این، نرم تصویر همه توزیع‌های احتمال برابر است:
\[
\|\Phi(P)\|_{\mathcal H}^2=\frac{\log2}{2}.
\]
پس تصاویر توزیع‌ها روی کره‌ای با شعاع ثابت
$r=\sqrt{\log2/2}$ در فضای هیلبرت قرار می‌گیرند.
بااین‌حال، \emph{این کره تمام فضای قابل‌دستیابی نیست}.
زیرا برای هر $a\ge0$ و $t>0$،
$\phi_a(t)\ge0$ است و همه مؤلفه‌های
$\Phi(P)(t)$ نامنفی‌اند. افزون بر آن،
مؤلفه‌ها باید هم‌زمان از یک بردار احتمال
با مجموع یک حاصل شده باشند. بنابراین:
\begin{equation}\label{cmp:eq:ch3-js-positive-spherical-image}
\boxed{
\begin{aligned}
\Phi(\Delta_K)\subseteq\bigl\{f\in\mathcal H:\;&
\|f\|_{\mathcal H}^{2}=\tfrac{\log2}{2},\\[-1pt]
& f_j(t)\ge0\
\text{برای تقریباً همه }t>0\bigr\}.
\end{aligned}
}
\end{equation}
این مجموعه، بخشی محدود از ناحیه نامنفی
کره هیلبرتی است، نه تمام آن؛
همچنین حتی همه نقاط نامنفی کره نیز
لزوماً تصویر توزیع احتمال نیستند.
تعبیر تصویریِ «بخشی شبیه یک هشتم کره»
تنها یک تشبیه برای محدودیت علامت مؤلفه‌هاست:
فضای هیلبرت حاضر بی‌نهایت‌بعدی است و
توزیع‌های روی $K$ برچسب فقط
$K-1$ درجه آزادی دارند.
از نامنفی‌بودن تصاویر،
$\langle\Phi(P),\Phi(Q)\rangle_{\mathcal H}\ge0$
و درنتیجه زاویه میان تصاویر
بیش از $\pi/2$ نخواهد بود.

\paragraph{از فاصلهٔ تصاویر به JSD و جذر آن.}
برای دو توزیع $P,Q\in\Delta_K$، مخلوط مشترک
$M=(P+Q)/2$ را در نظر می‌گیریم. واگرایی
جنسن--شنون بر حسب همین مخلوط تعریف می‌شود:
\begin{equation}\label{cmp:eq:ch3-JSD-definition}
\operatorname{JSD}(P,Q)
=\frac12D_{\mathrm{KL}}\!\left(P\middle\|M\right)
+\frac12D_{\mathrm{KL}}\!\left(Q\middle\|M\right),
\qquad M=\frac{P+Q}{2}.
\end{equation}
محاسبهٔ فاصلهٔ دو تصویر هیلبرتی با استفاده
از ضرب داخلی و نرم ثابت آن‌ها نشان می‌دهد:
\begin{equation}\label{cmp:eq:ch3-js-Hilbert-metric}
\boxed{
\|\Phi(P)-\Phi(Q)\|_{\mathcal H}^2
=\frac12D_{\mathrm{KL}}(P\|M)
+\frac12D_{\mathrm{KL}}(Q\|M)
=\operatorname{JSD}(P,Q),
\quad M=\frac{P+Q}{2}.
}
\end{equation}
بنابراین فاصله موردنظر به‌صورت زیر تعریف می‌شود:
\begin{equation}\label{cmp:eq:ch3-js-distance}
\boxed{
d_{\sqrt{\mathrm{JS}}}(P,Q)
=\|\Phi(P)-\Phi(Q)\|_{\mathcal H}
=\sqrt{\operatorname{JSD}(P,Q)}.
}
\end{equation}
این فاصله متقارن، یک‌به‌یک در تشخیص توزیع‌های متفاوت، کراندار در بازه $[0,\sqrt{\log2}]$ و واجد نامساوی مثلثی است؛ بنابراین یک متریک واقعی است. ساختار متریک جذر JSD و نمایش هیلبرتی آن در پژوهش‌های پیشین نیز بررسی شده‌اند \cite{EndresSchindelin2003,FugledeTopsoe2004}. این تعریف دوتوزیعی با وزن‌های برابر است؛
در تعمیم چندتوزیعی باید وزن‌ها و مخلوط مربوط
به‌صورت مستقل مشخص شوند.

\paragraph{خواص احتمالاتی و کران‌های JSD.}
از تعریف واگرایی جنسن--شنون می‌توان با
جایگذاری رابطهٔ KL، نمایش آنتروپی آن را نتیجه گرفت:
\begin{equation}\label{cmp:eq:ch3-JSD-entropy}
\operatorname{JSD}(P,Q)
=H\left(\frac{P+Q}{2}\right)
-\frac{H(P)+H(Q)}{2}.
\end{equation}
ازاین‌رو، JSD عدم‌تجانس دو توزیع
نسبت به مخلوط آن‌ها را اندازه می‌گیرد؛
با این تفاوت که برخلاف KL مستقیم،
به دلیل واردشدن همان مخلوط مشترک،
در صورت وجود مؤلفه‌های صفر نیز متناهی می‌ماند.
زیرا $M(j)\ge P(j)/2$ و $M(j)\ge Q(j)/2$،
هر یک از واگرایی‌های نسبت به مخلوط
حداکثر $\log2$ است. با نامنفی‌بودن KL
و شرط تساوی آن نتیجه می‌گیریم:
\[
0\le\operatorname{JSD}(P,Q)\le\log2,
\qquad
\operatorname{JSD}(P,Q)=0
\ \Longleftrightarrow\ P=Q.
\]
تقارن مستقیماً از تعریف به‌دست می‌آید؛
بیشینه $\log2$ در حالت تکیه‌گاه‌های
جدا به‌دست می‌آید. مثلاً برای
$P=(1,0)$ و $Q=(0,1)$،
KL مستقیم نامتناهی ولی JSD برابر $\log2$ است.
خود JSD عموماً نامساوی مثلثی را ارضا نمی‌کند،
اما جذر آن، مطابق نمایش هیلبرتی بالا، متریک است.


\paragraph{شباهت تصاویر و شباهت نرمال‌شده.}
از ثابت‌بودن نرم تصاویر و بسط مربع فاصله، رابطه زیر حاصل می‌شود:
\begin{equation}\label{cmp:eq:ch3-js-Hilbert-cosine}
\operatorname{JSD}(P,Q)
=\log2-2\langle\Phi(P),\Phi(Q)\rangle_{\mathcal H}.
\end{equation}
در نتیجه شباهت نرمال‌شده زیر دقیقاً برابر شباهت کسینوسی تصاویر است:
\begin{equation}\label{cmp:eq:ch3-js-similarity}
\boxed{
s(P,Q)=1-\frac{\operatorname{JSD}(P,Q)}{\log2}
=\frac{\langle\Phi(P),\Phi(Q)\rangle_{\mathcal H}}
{\|\Phi(P)\|_{\mathcal H}\|\Phi(Q)\|_{\mathcal H}},
\quad 0\le s(P,Q)\le1.
}
\end{equation}
این شباهت متقارن است؛ برای توزیع‌های یکسان مقدار یک
و برای تکیه‌گاه‌های جدا مقدار صفر دارد.
از آنجا که $s$ ضرب داخلی تصاویر هیلبرتی
نرمال‌شده است، یک \emph{کرنل مثبت نیمه‌معین}
محسوب می‌شود. به‌علاوه، JSD نسبت به هر توزیع
با ثابت‌بودن دیگری محدب است؛ بنابراین
$s=1-\operatorname{JSD}/\log2$
نسبت به هر ورودی به‌تنهایی مقعر خواهد بود.
برای مقایسه با بردارهای یک‌داغ نیز
شباهت $s(e_c,P)$ تابعی اکیداً صعودی
از مؤلفه $p_c$ است و ترتیب مؤلفه‌های
توزیع را حفظ می‌کند؛ البته این حکم خاص
مقایسه با بردارهای یک‌داغ است و درباره
مراکز عمومی به همان صورت ادعا نمی‌شود.
اثبات ویژگی‌های نگاشت، فرمول‌های انتگرالی،
کرنل مثبت نیمه‌معین، تقعر و خاصیت یک‌داغ
در بخش~\ref{app:probability-dissimilarity-proofs} آمده است.
پس متریک $d_{\sqrt{\mathrm{JS}}}$ و شباهت $s$
هر دو از \emph{همان تصاویر هیلبرتی با نرم ثابت}
به‌دست می‌آیند، اما نقش متفاوتی دارند:
اولی فاصله و دومی شباهت مبتنی بر ضرب داخلی است؛
$s$ خود متریک فاصله نیست.

\paragraph{دورترین توزیع از یک توزیع دلخواه.}
یکی از پیامدهای مهم این هندسه آن است
که برای هر توزیع ثابت $P\in\Delta_K$،
دست‌کم یک توزیع \emph{یک‌داغ}
بیشترین فاصله متریک جنسن--شنون را
از $P$ دارد. اگر $e_j$ رأس متناظر با
برچسب $j$ باشد، آنگاه
\begin{equation}\label{cmp:eq:ch3-js-farthest-vertex}
\boxed{
\begin{aligned}
\max_{Q\in\Delta_K}d_{\sqrt{\mathrm{JS}}}(P,Q)
&=\max_{1\le j\le K}d_{\sqrt{\mathrm{JS}}}(P,e_j)\\
&=d_{\sqrt{\mathrm{JS}}}(P,e_{j_*}),
\qquad j_*\in\arg\min_j p_j.
\end{aligned}}
\end{equation}
دلیل بخش نخست آن است که
$Q=\sum_jq_je_j$ و واگرایی
$\operatorname{JSD}(P,Q)$ نسبت به $Q$
\emph{محدب} است؛ پس
\[
\operatorname{JSD}(P,Q)
\le\sum_jq_j\operatorname{JSD}(P,e_j)
\le\max_j\operatorname{JSD}(P,e_j).
\]
برای تشخیص اینکه کدام رأس دورتر است،
از رابطه صعودی‌بودن
$s(P,e_j)$ نسبت به $p_j$ استفاده می‌کنیم:
\[
s(P,e_j)=
\frac{(1+p_j)\log(1+p_j)-p_j\log p_j}
{2\log2}.
\]
بنابراین کمترین مؤلفه $P$،
کمترین شباهت و بیشترین فاصله را
تولید می‌کند. در صورت برابر بودن
چند مؤلفه کمینه، چند رأس دورترین‌اند.
این گزاره \emph{وجود یک رأس بیشینه‌ساز}
را تضمین می‌کند، نه اینکه همه
بیشینه‌سازها حتماً رأس باشند:
مثلاً اگر $P=(1,0,0)$ باشد،
توزیع $Q=(0,\tfrac12,\tfrac12)$
نیز همان فاصله بیشینه $\sqrt{\log2}$
را دارد. اثبات شرایط محدب‌بودن و
مرز سیمپلکس در
بخش~\ref{app:probability-dissimilarity-proofs} آمده است.



\section{اثبات خواص نگاشت هیلبرتی و سنجه‌های جنسن--شنون}
\label{app:probability-dissimilarity-proofs}
در این بخش خواص معرفی‌شده در بخش~\ref{cmp:subsec:ch3-probability-distance-similarity} را اثبات می‌کنیم. نمادها و قرارداد مؤلفه‌های صفر همان قرارداد بخش پیشین‌اند.

\paragraph{تعریف نگاشت.}
فضای توزیع‌های احتمال روی $K$ برچسب را با
$\Delta_K=\{P\in\mathbb R_{\ge0}^{K}:\sum_{c=1}^{K}p_c=1\}$
نشان می‌دهیم. فضای هیلبرت
$\mathcal H=L^2((0,\infty);\mathbb R^K)$
را با ضرب داخلی معمول در نظر می‌گیریم و نگاشت
$\Phi:\Delta_K\longrightarrow\mathcal H$ را مستقیماً به‌صورت زیر تعریف می‌کنیم:
\begin{equation}\label{eq:app-ch3-js-map-definition}
\boxed{
\Phi(P)(t)=
\left(\frac{1-e^{-p_1t}}{2t},\ldots,
\frac{1-e^{-p_Kt}}{2t}\right),\qquad t>0.
}
\end{equation}
برای بررسی مؤلفه‌ها، فقط از نماد کوتاه
$\phi_a(t)=(1-e^{-at})/(2t)$، برای $a\ge0$، استفاده می‌کنیم. از
$1-e^{-at}\le\min\{at,1\}$ نتیجه می‌شود
$0\le\phi_a(t)\le\min\{a/2,1/(2t)\}$؛ بنابراین $\phi_a$ در صفر محدود و مربع آن در بی‌نهایت انتگرال‌پذیر است. در نتیجه $\Phi(P)$ برای هر $P\in\Delta_K$ عضوی از $\mathcal H$ است.

\paragraph{خواص نگاشت.}
برای $a\ge0$ و $t>0$، مشتق‌های مؤلفه نگاشت عبارت‌اند از:
\begin{equation}\label{eq:app-ch3-js-map-derivatives}
\begin{aligned}
\frac{\partial\phi_a(t)}{\partial a}
&=\frac12e^{-at}>0,\\
\frac{\partial^2\phi_a(t)}{\partial a^2}
&=-\frac t2e^{-at}<0,\\
\frac{\partial\phi_a(t)}{\partial t}
&=\frac{(at+1)e^{-at}-1}{2t^2}\le0.
\end{aligned}
\end{equation}
نامساوی آخر از $e^{at}\ge1+at$ ناشی می‌شود و برای $a>0$ اکید است. همچنین
\[
\lim_{t\to0^+}\phi_a(t)=\frac a2,\qquad
\lim_{t\to\infty}\phi_a(t)=0.
\]
از مشتق اول، برای هر $t>0$ داریم:
\[
a>b\ \Longleftrightarrow\ \phi_a(t)>\phi_b(t).
\]
پس نگاشت ترتیب احتمال‌ها را میان مؤلفه‌های تصویر حفظ می‌کند. همچنین اگر
$\Phi(P)=\Phi(Q)$ در $\mathcal H$ باشد، مؤلفه‌های دو تصویر تقریباً همه‌جا برابرند؛ به دلیل پیوستگی آن‌ها در $t>0$، این برابری در همه $t>0$ برقرار است. از اکیداً صعودی بودن $\phi_a(t)$ نسبت به $a$ نتیجه می‌شود $p_c=q_c$ برای تمام $c$. بنابراین $\Phi$ یک‌به‌یک است.

\paragraph{محاسبه ضرب داخلی و نرم تصویرها.}
مستقیماً از تعریف نگاشت، برای $a,b>0$ داریم:
\[
\langle\phi_a,\phi_b\rangle_{L^2}
=\frac14\int_0^\infty
\frac{(1-e^{-at})(1-e^{-bt})}{t^2}\,dt.
\]
عبارت داخل انتگرال پس از انتگرال‌گیری جزءبه‌جزء به‌صورت زیر درمی‌آید:
\begin{equation}\label{eq:app-ch3-js-integral-representation}
\begin{aligned}
4\langle\phi_a,\phi_b\rangle_{L^2}
&=\int_0^\infty\frac{
a(e^{-at}-e^{-(a+b)t})
+b(e^{-bt}-e^{-(a+b)t})}{t}\,dt\\
&=a\log\frac{a+b}{a}
+b\log\frac{a+b}{b}.
\end{aligned}
\end{equation}
برای تساوی دوم از هویت فرولانی
$\int_0^\infty(e^{-ut}-e^{-vt})\,dt/t=\log(v/u)$
برای $u,v>0$ استفاده شده است. عبارت مرزی انتگرال‌گیری جزءبه‌جزء در هر دو انتهای بازه صفر است. رابطه با قرارداد حدی $0\log0=0$ برای $a=0$ یا $b=0$ نیز معتبر است و در حالت $a=b=0$ هر دو طرف صفرند.

حال با قرار دادن $a=b$، داریم:
\[
\|\phi_a\|_{L^2}^2=\frac{a\log2}{2}.
\]
درنتیجه، به دلیل $\sum_cp_c=1$، نرم تصویر هر توزیع ثابت است:
\begin{equation}\label{eq:app-ch3-js-sphere}
\boxed{
\|\Phi(P)\|_{\mathcal H}^2
=\sum_{c=1}^{K}\|\phi_{p_c}\|_{L^2}^2
=\frac{\log2}{2}.
}
\end{equation}
بنابراین تمام تصاویر توزیع‌های احتمال روی کره‌ای به شعاع
$\sqrt{\log2/2}$ در $\mathcal H$ قرار می‌گیرند.

\paragraph{محدودبودن تصویر به ناحیه مجاز کره.}
برای $P\in\Delta_K$، از
$\phi_{p_j}(t)=(1-e^{-p_jt})/(2t)\ge0$
نتیجه می‌شود تمام مؤلفه‌های
$\Phi(P)(t)$ نامنفی هستند. بنابراین
تصویر $\Phi(\Delta_K)$ زیرمجموعه تقاطع
کره‌ای با شعاع $\sqrt{\log2/2}$
و مخروط توابع نامنفی $\mathcal H$ است.
\emph{عکس این حکم درست نیست}:
هر تابع دلخواه نامنفی روی این کره،
لزوماً به شکل $\Phi(P)$ نیست،
زیرا باید برای تمامی $j$ و تمامی
$t>0$ از یک مجموعه ثابت ضرایب
$p_j\ge0$ با مجموع یک حاصل شود.
این تصویر تنها $K-1$ درجه آزادی
در فضای هیلبرت بی‌نهایت‌بعدی دارد.
به‌علاوه، برای هر دو توزیع
$P,Q$، ضرب داخلی تصاویر آن‌ها
نامنفی است؛ بنابراین زاویه هیلبرتی
این تصاویر در بازه $[0,\pi/2]$ قرار دارد.
در مقابل، تصاویر بردارهای یک‌داغ
$e_j$ و $e_k$ برای $j\ne k$
متعامدند: برای این دو توزیع
تکیه‌گاه جدا، $\operatorname{JSD}(e_j,e_k)=\log2$
و از رابطه~\eqref{cmp:eq:ch3-js-Hilbert-cosine}
ضرب داخلی صفر می‌شود.
پس عبارت «بخشی از کره» بیان
دقیقی از هندسه این تصویر است،
و «یک هشتم کره» فقط تشبیهی
غیردقیق به یک ناحیه نامنفی
در فضای سه‌بعدی خواهد بود.

\paragraph{فاصله دو تصویر و رابطه آن با جذر JSD.}
با بسط مربع فاصله و جایگذاری رابطه ضرب داخلی:
\begin{equation}\label{eq:app-ch3-js-Hilbert-distance}
\begin{aligned}
\|\Phi(P)-\Phi(Q)\|_{\mathcal H}^2
&=\|\Phi(P)\|^2+\|\Phi(Q)\|^2
-2\langle\Phi(P),\Phi(Q)\rangle\\
&=\log2-\frac12\sum_{c=1}^{K}
\left[
p_c\log\frac{p_c+q_c}{p_c}
+q_c\log\frac{p_c+q_c}{q_c}
\right]\\
&=\frac12\sum_{c=1}^{K}\left[
p_c\log\frac{2p_c}{p_c+q_c}
+q_c\log\frac{2q_c}{p_c+q_c}
\right]\\
&=\operatorname{JSD}(P,Q).
\end{aligned}
\end{equation}
در تمام سطرها، عبارت‌های دارای ضریب صفر با مقدار حدی آن‌ها تفسیر می‌شوند. بنابراین:
\[
\boxed{
\|\Phi(P)-\Phi(Q)\|_{\mathcal H}
=\sqrt{\operatorname{JSD}(P,Q)}.
}
\]
این برابری نشان می‌دهد که نگاشت $\Phi$ \emph{ایزومتریک} است؛ به‌ویژه از یک‌به‌یک بودن نگاشت و خواص نرم هیلبرتی نتیجه می‌شود جذر JSD یک متریک است و:
\begin{equation}\label{eq:app-ch3-js-triangle}
d_{\sqrt{\mathrm{JS}}}(P,R)\le
d_{\sqrt{\mathrm{JS}}}(P,Q)+d_{\sqrt{\mathrm{JS}}}(Q,R).
\end{equation}
به‌علاوه، چون همه مؤلفه‌های $\Phi(P)(t)$ و $\Phi(Q)(t)$ نامنفی‌اند، ضرب داخلی آن‌ها نامنفی است. از رابطه مربع فاصله نتیجه می‌شود
$0\le\operatorname{JSD}(P,Q)\le\log2$؛
پس فاصله مزبور متناهی و کراندار است.

\paragraph{شباهت تصاویر و رابطه آن با شباهت نرمال‌شده.}
از محاسبه فاصله دو تصویر و ثابت‌بودن نرم‌ها داریم:
\[
2\langle\Phi(P),\Phi(Q)\rangle_{\mathcal H}
=\log2-\operatorname{JSD}(P,Q).
\]
پس با تعریف شباهت کسینوسی دو تصویر:
\begin{equation}\label{eq:app-ch3-js-cosine}
\begin{gathered}
\frac{\langle\Phi(P),\Phi(Q)\rangle_{\mathcal H}}
{\|\Phi(P)\|_{\mathcal H}\|\Phi(Q)\|_{\mathcal H}}\\
=\frac{2\langle\Phi(P),\Phi(Q)\rangle_{\mathcal H}}{\log2}\\
=1-\frac{\operatorname{JSD}(P,Q)}{\log2}
=s(P,Q).
\end{gathered}
\end{equation}
بنابراین $d_{\sqrt{\mathrm{JS}}}$ فاصلهٔ هیلبرتی تصاویر است؛
$s$ نیز شباهت کسینوسی همان تصاویر را نشان می‌دهد.
هر دو متقارن‌اند و $0\le s(P,Q)\le1$.
همچنین $s(P,Q)=1$ اگر و فقط اگر $P=Q$.

از نامنفی بودن $\phi_a(t)$
و مثبت بودن اکید آن برای $a>0$ نتیجه می‌شود
که $s(P,Q)=0$ دقیقاً در صورت جدا بودن
تکیه‌گاه‌های $P$ و $Q$ رخ می‌دهد.
افزون بر این، $s$ یک کرنل مثبت نیمه‌معین است.
برای ضرایب حقیقی $\alpha_1,\ldots,\alpha_n$ داریم:
\[
\begin{aligned}
\sum_{i,j=1}^{n}\alpha_i\alpha_j s(P_i,P_j)
&=\frac{2}{\log2}
\left\|\sum_{i=1}^{n}\alpha_i\Phi(P_i)\right\|_{\mathcal H}^2\\
&\ge0.
\end{aligned}
\]
در مقابل، $s$ خود متریک فاصله نیست؛ دو معیار از ساختار هندسی یکسانی به‌دست می‌آیند، اما نقش‌های متفاوتی دارند.

\paragraph{حفظ ترتیب احتمال‌ها در امتیاز طبقه‌ها.}
یک توزیع $P=(p_1,\ldots,p_K)$ و بردار یک‌داغ $e_c$ را در نظر بگیریم. چون تنها مؤلفه $c$ در تصویر $e_c$ ناصفر است، رابطه ضرب داخلی به‌صورت زیر ساده می‌شود:
\[
\langle\Phi(e_c),\Phi(P)\rangle_{\mathcal H}
=\frac14\left[
\log(1+p_c)+p_c\log\frac{1+p_c}{p_c}
\right].
\]
با اعمال نرمال‌سازی کسینوسی نتیجه می‌شود:
\begin{equation}\label{eq:app-ch3-js-onehot-monotonic}
\begin{aligned}
s(e_c,P)
&=\frac{(1+p_c)\log(1+p_c)-p_c\log p_c}{2\log2},\\
\frac{d\,s(e_c,P)}{dp_c}
&=\frac{1}{2\log2}\log\frac{1+p_c}{p_c}>0,
\qquad 0<p_c<1.
\end{aligned}
\end{equation}
به دلیل پیوستگی تابع در $p_c=0$ و $p_c=1$، این شباهت در کل $[0,1]$ اکیداً صعودی است؛ بنابراین:
\[
p_c>p_r\quad\Longleftrightarrow\quad
s(e_c,P)>s(e_r,P).
\]
پس انتخاب بیشترین شباهت با بردارهای یک‌داغ،
ترتیب احتمالات را حفظ می‌کند.
از طرف دیگر،
\[
d_{\sqrt{\mathrm{JS}}}(e_c,P)^2
=\log2\,[1-s(e_c,P)].
\]
بنابراین انتخاب کمترین فاصله نیز همین ترتیب را دارد.
این خاصیت به مقایسه با بردارهای یک‌داغ اختصاص دارد.
در طبقه‌بند واقعی، توزیع‌های همسایگی
با مراکز طبقه‌ها مقایسه می‌شوند؛
برابری با احتمال پسین واقعی ادعا نمی‌شود.

\paragraph{تقعر شباهت نسبت به توزیع‌ها.}
محاسبه مشتق دوم از فرمول JSD به‌دست‌آمده از فاصله نشان می‌دهد که با ثابت گرفتن $Q$ و برای $p_c>0$:
\[
\frac{\partial^2\operatorname{JSD}(P,Q)}{\partial p_c^2}
=\frac{q_c}{2p_c(p_c+q_c)}\ge0,
\qquad
\frac{\partial^2\operatorname{JSD}(P,Q)}
{\partial p_c\partial p_r}=0\quad(c\ne r).
\]
بنابراین JSD نسبت به $P$ محدب و $s(P,Q)$ نسبت به $P$ مقعر است. به دلیل تقارن، حکم مشابه برای $Q$ با ثابت‌بودن $P$ نیز برقرار است. نتیجه به مرز سیمپلکس با پیوستگی تعمیم می‌یابد.

\paragraph{بیشینه فاصله JS از یک توزیع ثابت.}
می‌خواهیم نشان دهیم برای هر
$P=(p_1,\ldots,p_K)\in\Delta_K$،
بیشینه فاصله تا $Q\in\Delta_K$
در دست‌کم یکی از رأس‌های $e_j$
سیمپلکس به‌دست می‌آید.
از رابطه مشتق دوم بالا و تقارن JSD،
$Q\mapsto\operatorname{JSD}(P,Q)$
روی درون سیمپلکس محدب است
و این خاصیت با پیوستگی
تا مرز سیمپلکس ادامه می‌یابد.
هر $Q$ را به شکل ترکیب محدب
$Q=\sum_{j=1}^{K}q_je_j$ می‌نویسیم.
نامساوی ینسن می‌دهد:
\begin{equation}\label{eq:app-ch3-js-farthest-vertex-convexity}
\begin{aligned}
\operatorname{JSD}(P,Q)
&=\operatorname{JSD}\!\left(
P,\sum_jq_je_j\right)\\
&\le\sum_jq_j\operatorname{JSD}(P,e_j)\\
&\le\max_j\operatorname{JSD}(P,e_j).
\end{aligned}
\end{equation}
از طرف دیگر، خود رأس‌ها
عضو $\Delta_K$ هستند؛
پس این کران قابل‌دستیابی است و
بیشینه روی رأس‌ها تحقق می‌یابد.
برای مشخص‌کردن رأس دورتر،
رابطه~\eqref{eq:app-ch3-js-onehot-monotonic}
به ما می‌دهد:
\begin{equation}\label{eq:app-ch3-js-farthest-vertex-formula}
\begin{aligned}
\operatorname{JSD}(P,e_j)
&=\log2\,[1-s(P,e_j)]\\
&=\log2-\frac12\left[
(1+p_j)\log(1+p_j)-p_j\log p_j
\right].
\end{aligned}
\end{equation}
با قرارداد $0\log0=0$ و
پیوستگی در کرانه‌ها، تابع
$g(p)=(1+p)\log(1+p)-p\log p$
روی $[0,1]$ اکیداً صعودی است؛
زیرا برای $0<p<1$ داریم
$g'(p)=\log((1+p)/p)>0$.
پس $j_*\in\arg\min_jp_j$
رأسی با \emph{کمترین شباهت}
و \emph{بیشترین فاصله} خواهد بود.
اگر $p_*=\min_jp_j$ باشد،
مقدار دقیق بیشینه عبارت است از:
\begin{equation}\label{eq:app-ch3-js-max-distance-value}
\boxed{
\max_{Q\in\Delta_K}d_{\sqrt{\mathrm{JS}}}(P,Q)^2
=\log2-\frac12\left[
(1+p_*)\log(1+p_*)-p_*\log p_*
\right].
}
\end{equation}
در نتیجه اگر همه مؤلفه‌های $P$
مثبت باشند، بیشترین فاصله
از کران عمومی $\sqrt{\log2}$
\emph{کوچک‌تر} است.
اگر دست‌کم یک مؤلفه صفر باشد،
بیشترین فاصله برابر
$\sqrt{\log2}$ خواهد بود.
در حالت اخیر، بیشینه‌سازهای غیررأسی نیز
ممکن‌اند. هر $Q$ با تکیه‌گاه محدود به
مؤلفه‌های صفر $P$، تکیه‌گاهی جدا از $P$
و فاصلهٔ بیشینه دارد.
برای نمونه، در $K=3$ داریم:
\[
P=(1,0,0),\qquad
Q=\left(0,\frac12,\frac12\right).
\]
این مثال یک بیشینه‌ساز غیررأسی است.
پس نتیجهٔ اثبات‌شده، \emph{وجود دست‌کم
یک رأس دورترین} است؛
نه اینکه هر بیشینه‌ساز لزوماً
برداری یک‌داغ باشد.

